\chapter{绪论}\label{chap:introduction}

\section{研究背景与意义}

多功能雷达网络在电子对抗任务中通常需要同时回答两个问题：目标区域是否被可靠探测，重点任务点是否受到足够的压制。单部雷达很难在复杂电磁环境下兼顾这两类需求，因此工程上更常采用多节点协同组网。已有综述表明，组网性能除单站体制和信号处理能力外，还受节点位置、融合结构和抗干扰条件共同影响~\cite{Javadi2020}。本文讨论的部署优化，正是把这些因素中的“节点位置”抽出来进行定量设计。

空中平台的加入使部署问题更复杂。无人机集群研究表明，复杂任务中的协同飞行、通信可靠性、异构平台分工和环境扰动会共同影响任务规划~\cite{Alqudsi2025UAVSwarms}。在防空与反无人机场景中，多类型雷达需要在复杂地形下形成覆盖互补，近期研究已把地形遮挡、优先探测区域和多类型雷达协同纳入部署优化~\cite{Ding2025UAVRadarDeployment}。对本文的空地协同问题而言，空中节点通常拥有更好的视距传播条件，适合弥补地面节点在遮挡区域中的覆盖不足；地面节点则便于持续驻留，可为局部区域提供较稳定的干扰压制。若节点数量有限，优化器必须在“扩大覆盖面”和“抬高最薄弱点压制强度”之间做取舍，这个取舍不适合用单一指标排序，而应以Pareto前沿形式给出一组候选方案。

另一个直接来自场景几何的困难是边界效应。部署区域可能是非凸多边形，也可能包含空洞、禁入区或多个分离子区域。如果先在外接矩形内采样，再删除落在区域外的点，靠近凹边界和窄通道的可行位置会被明显削弱；这些位置恰好可能影响边缘任务点的覆盖。与此同时，若在空地协同场景中仍使用同一套路径损耗参数，空中链路与地面链路的差异会被折叠掉，目标函数得到的部署优劣也会偏离实际传播条件。

已有MOPSO-DT方法通过凸分解和坐标变换处理复杂部署区域，为本文提供了可复用的几何框架~\cite{Han2025}。本文在此基础上做两点延伸：一是把空地链路和地地链路分开建模，二是在相同编码、相同目标函数和相同评价预算下加入NSGA-II、MOEA/D、SPEA2等基线，避免只在同源MOPSO变体之间比较。图~\ref{fig:intro_research_framework}概括了本文从任务需求、区域约束、传播建模到Pareto方案输出的研究主线。本文关注的不是证明某个优化器在所有指标上都占优，而是评估复杂区域约束、异构传播模型和参数策略放在同一流程中后能带来哪些可观察变化。

\begin{figure}[!htbp]
\centering
\includegraphics[width=0.88\linewidth]{figures/intro_research_framework.pdf}
\caption{本文研究问题与技术路线。}
\label{fig:intro_research_framework}
\end{figure}

\section{国内外研究现状}

雷达网络部署优化方面，早期工作常把覆盖率、成本或站址可行性作为单一目标，用遗传算法、模拟退火等元启发式方法搜索节点位置。随着任务指标增多，研究逐渐转向多目标表述，包括组网雷达搜索与跟踪联合部署、区域约束下的混合变量编码，以及非连通区域部署优化等问题~\cite{Yan2024,Han2025,MSRS2024}。面向多类型雷达反无人机覆盖的研究把复杂地形、优先探测区域和算法稳定性作为部署质量的核心因素~\cite{Ding2025UAVRadarDeployment}；面向分布式无源雷达的几何优化则开始同时考虑节点位置、阵列指向和监视覆盖~\cite{Tian2026PassiveRadarGeo}。面向干扰源协同部署和分布式雷达拓扑的研究也显示，节点位置与任务区域、探测指标、干扰指标之间存在强耦合~\cite{JammerDeploy2024,Cao2026NSMOPSO}。这些工作说明多目标建模是必要的，但边界覆盖和空地异构传播往往没有被放在同一实验框架下检验。

复杂区域约束与边界效应处理方面，非凸边界、空洞和禁入区会使连续优化器频繁产生非法候选点。计算几何中的凸分解为这类区域提供了直接处理方式，相关研究已经覆盖快速多边形分解、几何查询和带空洞多边形切割等基础问题~\cite{Hertel1983,DeBerg2008,Chazelle1982}。近年的Radar3D研究则把地形和传播环境纳入雷达部署建模~\cite{Radar3D2025}。现有方法通常能够保证候选点可行，但对“边界任务点是否因此受益”的验证不够充分。本文单独设置边界ECR实验，目的就是把这一几何处理的效果从总体覆盖率中拆出来观察。

多目标优化算法方面，进化多目标优化的代表性路线包括非支配排序、分解式求解、外部档案和密度估计等~\cite{Vachhani2015,Hua2021}。NSGA-II、SPEA2和MOEA/D分别通过快速非支配排序、强度Pareto适应度和分解式标量子问题维护非支配解集~\cite{Deb2002,Zitzler2001,Zhang2007}。PSO路线的优势在于连续空间更新形式简单，MOPSO基本框架和拥挤距离领导者选择为本文档案维护提供了直接依据~\cite{Coello2004,Raquel2005}。近年多种MOPSO改进方法又从区域选择、领导者预测、指标融合、动态小生态和分解选择等角度改进搜索过程~\cite{Liu2023,Wang2024,Zhang2024,Xu2024,Sun2024,Ma2025}。对雷达部署而言，算法名称并不是结论本身；更需要控制编码、评价函数和预算，使不同优化器在同一问题上比较。

空地异构传播建模方面，空地协同部署区别于纯地面网络的地方在于节点高度和传播路径。航空移动业务、地面站间传播和通用路径损耗模型为A2G/G2G链路区分提供了基础依据~\cite{ITUP528,ITUP452,Rappaport2002}。UAV协同、ISAC部署和城市空中交通监视等研究进一步推动了空中平台部署与覆盖优化结合，不同场景下的空中节点会改变覆盖设计方式~\cite{Li2025,Milone2024}；UAV集群综述也指出，异构平台、通信协同和任务分配是未来复杂无人系统的重要研究方向~\cite{Alqudsi2025UAVSwarms}。若把所有节点都视作同质地面站，模型实现会更简单，但会丢失空中节点在视距传播和覆盖半径上的差异。本文因此在目标函数计算中显式区分A2G和G2G链路。

综上，已有研究分别提供了雷达部署、多目标优化、复杂几何处理和传播建模的基础。本文把这些环节组合为一个可复现实验流程：先用凸分解和坐标变换处理复杂区域，再用异构传播模型计算覆盖和干扰目标，最后在同一预算下比较MOPSO-DT改进配置与若干经典基线。这样得到的结论更容易对应到具体组件，而不是停留在单一算法性能描述上。

\section{研究内容与创新点}

本文围绕三个问题展开：复杂区域中的候选点怎样保持可行且不系统性远离边界；空地协同网络中的传播差异怎样进入目标函数；覆盖率ECR与最小等效干扰强度$J_{\min}$发生冲突时，如何给出可供任务选择的非支配方案。

本文主要创新点如下：

\begin{enumerate}
    \item 面向复杂区域的可行部署编码。本文使用Hertel-Mehlhorn凸分解和垂直交点坐标变换，将非凸、带空洞或不连通区域转化为可编码凸子区域及其内部连续坐标。该处理减少了外接矩形采样带来的非法点过滤，并通过边界ECR实验单独检查边缘任务点的覆盖变化。
    \item 区分A2G/G2G链路的目标函数。在覆盖和干扰评估中，本文分别设置空地链路与地地链路的路径损耗指数，并以ECR和$J_{\min}$构造双目标最小化问题。这样可以保留空中节点的传播优势，也能观察其对最薄弱任务点压制强度的影响。
    \item 同预算多算法对比与权衡分析。本文在同一混合变量编码、同一评价函数和同一评价预算下比较改进MOPSO-DT、原始MOPSO-DT、NSGA-II、MOEA/D、SPEA2和随机搜索，并结合参数调优、消融实验和边界专项实验解释各组件的作用。实验结果用于说明部署流程的适用边界，而不是泛化为所有场景下的算法排序。
\end{enumerate}

\section{论文结构安排}

全文共六章。第一章说明问题来源、相关研究和本文工作范围。第二章给出多目标优化、粒子群优化和计算几何基础。第三章定义空地协同部署场景、探测/干扰模型和双目标优化问题。第四章描述区域分解、坐标变换、MOPSO-DT更新规则及公平对比协议。第五章报告多算法对比、边界效应、参数调优和消融实验。第六章总结本文结论，并说明实验规模、传播模型和统计检验方面的限制。
